\section{A compact positive density extension}\label{sqt:sec:density}
The square proof extends with uniform constants on each fixed compact interval of positive densities. This section specifies the quantifiers and all changes to the analytic argument; it does not assert uniformity in sparse or unbounded density regimes.

Let $T(n,s)$ count labelled $n\times n$ nonnegative integer tables with all line sums $s$, and write
\begin{equation}\label{sqt:eq:densitydefs}
\lambda=\frac{s}{n},\qquad u=\lambda(1+\lambda),\qquad A_\lambda=\frac u2,
\qquad H(\lambda)=\frac{(1+\lambda)^{1+\lambda}}{\lambda^\lambda}.
\end{equation}
The matrix $A$ of \eqref{sqt:eq:augmented} is unrelated to the scalar $A_\lambda$; the whitening map $T$ stays as in \eqref{sqt:eq:whitening}. Define
\begin{equation}\label{sqt:eq:densityG}
\mathcal G(n,\lambda)=\frac{H(\lambda)^{n^2}}
 {(2\pi u)^{n-1/2}n^{n-1}},\qquad
B_0(u)=\frac13-\frac1{6u}.
\end{equation}

\begin{theorem}\label{sqt:thm:density}
There are polynomials $P_j\in\mathbb Q[z]$, of degree at most $j+1$, such that for every fixed compact interval $I=[\lambda_-,\lambda_+]\subset(0,\infty)$ and every fixed integer $K\geq0$,
\begin{equation}\label{sqt:eq:densitytheorem}
\log\frac{T(n,s)}{\mathcal G(n,s/n)}
 =B_0(u)+\sum_{j=1}^K\frac{P_j(1/u)}{n^j}
   +O_{I,K}(n^{-K-1}),
\end{equation}
uniformly over positive integers $n,s$ with $s/n\in I$. In particular,
\begin{equation}\label{sqt:eq:densitypolys}
P_1(z)=-\frac32,\qquad
P_2(z)=\frac{1171}{180}+\frac{11}{12}z+\frac1{60}z^2+\frac1{180}z^3.
\end{equation}
At $\lambda=1$, this recovers Theorem~\ref{sqt:thm:main}.
\end{theorem}
The statement permits varying admissible densities $s_n/n$. For a fixed rational density, it applies along the admissible subsequence. At a fixed irrational density, exact integral margins are unavailable; rounding $\lambda n$ changes the density and must be reflected in every occurrence of $\lambda$ in \eqref{sqt:eq:densitytheorem}.

\subsection{Uniform localization and the scaled gauge}
The centered characteristic function is now
\[
\psi_\lambda(z)=\frac{e^{-\ii\lambda z}}{1-\lambda(e^{\ii z}-1)},\qquad
\mathcal F_\lambda(\theta,\phi)=\prod_{i,j}\psi_\lambda(\theta_i+\phi_j).
\]
Although an individual factor need not be $2\pi$-periodic, the full product is periodic in each angle because $n\lambda=s\in\mathbb Z$. Gauge invariance is exact. Each table has geometric product probability $H(\lambda)^{-n^2}$, giving
\[
T(n,s)=H(\lambda)^{n^2}(2\pi)^{-2n}\int\mathcal F_\lambda.
\]

Canfield--McKay's condition~(1.1) specializes to
\[
\frac23\left(4+\frac1u\right)\leq a\log n.
\]
With fixed $a=b=1/8$, it holds uniformly on $I$ for all sufficiently large $n$; their polynomial density bound is automatic. Their stated asymptotic constants allow dependence on $a,b$ and the fixed localization exponent, with margins among the asymptotic variables. One can also check uniformity directly in their localization proof. The subdivision count $\lceil6000(1+\lambda)\rceil$ is bounded on $I$, its angular width is bounded above and below, and $A_\lambda$, $A_\lambda^{-1}$ and $\lambda/(1+\lambda)$ have the needed compact bounds. The far-region decay is $\exp\{-c_I n^{1+\eta}\}$ against a reciprocal Gaussian normalization $\exp\{O_I(n\log n)\}$. The exceptional-coordinate decay is $\exp\{-c_I n(m_2+n_2)\}$, and the remaining omitted small intervals contribute $\exp\{-c_I n^{2\eta}\}$. Their Lemma~4 is uniform in $\lambda>0$. More precisely, after normalization the three bounds have the forms
\[
 e^{C_I n\log n-c_I n^{1+\eta}},\qquad
 \sum_{1\leq r\leq2n^\eta}n^{C_Ir+C_I}e^{-c_Inr},\qquad
 n^{C_I}e^{-c_In^{2\eta}},
\]
where $r=m_2+n_2$ counts exceptional coordinates. The middle sum decays geometrically in $r$ for sufficiently large $n$. Each expression is $O_I(e^{-c'_In^\eta})$ for a sufficiently small fixed $\eta$. Thus
\begin{equation}\label{sqt:eq:densityminor}
\int_{\mathcal R_{\lambda,\eta}^c}|\mathcal F_\lambda|
 =O_I(e^{-c_I n^\eta})I_{0,\lambda},\qquad
I_{0,\lambda}=\pi^{n-1/2}A_\lambda^{-n+1/2}n^{1-n}
 e^{-(1+2A_\lambda)/(12A_\lambda)}.
\end{equation}
The region has centered deviations bounded by $(1+\lambda)^{-1}n^{-1/2+\eta}$ and mean sum bounded by $(1+\lambda)^{-1}n^{-1+4\eta}$. These bounds make its small $v^\perp$ representative $o_I(\rho)$ uniformly, as before.

The correct augmented weight is $\exp\{-A_\lambda(v\cdot x)^2/2\}$. Its quadratic part is $-A_\lambda q_+(x)$ and its full Gaussian integral is
\[
Z_{n,\lambda}=\frac1{\sqrt2}\left(\frac\pi{A_\lambda n}\right)^n.
\]
The real gauge fiber has mass $\sqrt{2\pi/A_\lambda}/(2n)$, so the conversion factor is $Q_{n,\lambda}=2n\sqrt{2\pi A_\lambda}$. In particular,
\[
\frac{Q_{n,\lambda}Z_{n,\lambda}}{(2\pi)^{2n}}
 =(4\pi A_\lambda)^{-n+1/2}n^{1-n},\qquad
\frac{I_{0,\lambda}}{Q_{n,\lambda}Z_{n,\lambda}}
 =\frac{e^{-(1+2A_\lambda)/(12A_\lambda)}}{2\pi}.
\]
The last ratio is uniformly bounded above and below on $I$. The proof of Lemma~\ref{sqt:lem:transfer} now applies verbatim with this weight: the omitted gauge tail is $e^{-\Omega_I(n^{1+2\eps})}$ and its normalization costs at most $e^{O_I(n\log n)}$. This proves a uniform torus-to-cube transfer. Using an unscaled gauge while retaining the scaled Gaussian covariance would give an incorrect normalization.

\subsection{Uniform analytic and cumulant estimates}
Define
\[
h_\lambda(z)=-\ii\lambda z-\log[1-\lambda(e^{\ii z}-1)]+A_\lambda z^2
 =\sum_{d\geq3}h_d(\lambda)z^d.
\]
The nearest singularity is at $-\ii\log(1+1/\lambda)$. For example, $r_0=\tfrac14\log(1+1/\lambda_+)>0$ gives a uniform analytic disk of radius $2r_0$, because
$\lambda_+(e^{2r_0}-1)<1$. The logarithmic series on this disk is uniformly bounded. Hence \eqref{sqt:eq:derivativebounds} holds with constants depending only on $I$, and
\begin{equation}\label{sqt:eq:densitycumulants}
h_d(\lambda)=\frac{\ii^d}{d!}\kappa_d(\lambda),\qquad
\kappa_d(\lambda)=\sum_{j=1}^d\genfrac{\{}{\}}{0pt}{}{d}{j}(j-1)!\lambda^j.
\end{equation}
The independent Gaussian coordinates before applying $T$ now have variance $1/(2A_\lambda n)$. Conditioning on the same cube preserves independence, with complement probability $O_I(e^{-c_I n^{2\eps}})$. The coefficient-vector bounds \eqref{sqt:eq:rowsumbounds} are unchanged. Therefore \eqref{sqt:eq:Dk} holds uniformly, and the same $\alpha,m,L$ in \eqref{sqt:eq:cutofforders} satisfy Isaev's hypotheses uniformly for sufficiently large $n$.

Taylor replacement inside the finite cumulants has the same exponent \eqref{sqt:eq:replacementexponent}; its constants depend on $I,p$. Fixed-degree Gaussian moments remain uniformly polynomially bounded since $A_\lambda$ stays away from zero. Removing the conditioning again costs less than every inverse power of $n$. The resulting finite exponent is real and uniformly bounded by the connected-diagram estimate below, so absolute localization errors are relative negligible errors. This establishes the uniform version of \eqref{sqt:eq:analyticreduction}.

\subsection{The density polynomials and their coefficients}
The pair-sum covariance becomes
\[
\Cov(Z_{ij},Z_{k\ell})=\frac{n\delta_{ik}+n\delta_{j\ell}-1}{2A_\lambda n^2}.
\]
For a tuple of even total degree $2E$, its Gaussian cumulant is exactly $A_\lambda^{-E}K_{\boldsymbol d}(n)$, with $K_{\boldsymbol d}$ the density-one polynomial. Hence the uniform finite formula is
\begin{equation}\label{sqt:eq:densityfinite}
\log\frac{T(n,s)}{\mathcal G(n,\lambda)}
 =\sum_{\substack{\boldsymbol d\text{ nondecreasing},\ r\geq1,\ d_a\geq3\\
                    D\text{ even},\ c(\boldsymbol d)\leq p}}
 \frac{\prod_a h_{d_a}(\lambda)}{\prod_d m_d!}
 A_\lambda^{-E}K_{\boldsymbol d}(n)+O_{I,p}(n^{-p}).
\end{equation}
The power bound is unchanged. This already proves existence, uniformity and a finite rational-function coefficient algorithm.

To obtain the polynomial structure, use
\[
\kappa_2=u,\qquad \kappa_{d+1}=u\frac{\dd}{\dd\lambda}\kappa_d.
\]
Inductively, $\kappa_d$ is divisible by $u$, has degree $d$, and satisfies
$\kappa_d(-1-\lambda)=(-1)^d\kappa_d(\lambda)$. For an even-total-degree tuple,
$\prod_a\kappa_{d_a}/u^r$ is a polynomial invariant under $\lambda\mapsto-1-\lambda$ of degree at most $2(E-r)$. Any such polynomial is a polynomial in $u$ of degree at most $E-r$: set $w=1+2\lambda$ and note that invariance means an even polynomial in $w$, with $w^2=1+4u$. Multiplication by $A_\lambda^{-E}=(2/u)^E$ therefore makes each weighted tuple a polynomial in $1/u$ of degree at most $E-r=c$. To contribute to $n^{-j}$ it must have $c\leq j+1$, proving the asserted degree bound.

The cost-one constant is
\[
\frac{1+6u}{2u}-\frac{2(1+4u)}{3u}=B_0(u).
\]
For the first correction, put $w=1+2\lambda$. The seven tuple contributions, in the order $(3,3),(4),(3,3,3,3),(3,3,4),(3,5),(4,4),(6)$, are
\begin{align*}
&\frac{w^2}{2u},\quad-\frac{1+6u}{2u},\quad
\frac{11w^4}{8u^2},\quad-\frac{31w^2(1+6u)}{12u^2},\\
&\frac{5w^2(1+12u)}{6u^2},\quad
\frac{13(1+6u)^2}{24u^2},\quad-\frac{1+30u+120u^2}{6u^2}.
\end{align*}
With $w^2=1+4u$, the first two sum to $-1$ and the remaining five to $-1/2$. All eighteen tuples at costs at most three give
\[
P_2(1/u)=\frac{1171u^3+165u^2+3u+1}{180u^3}.
\]
The same independent formal-factorization computation described in Section~\ref{sqt:sec:coefficients}, with a separately expanded geometric logarithm and the scalar factor $A_\lambda^{-E}$, verifies these rational-function identities exactly. Taking $p=K+1$ in \eqref{sqt:eq:densityfinite} completes the proof of Theorem~\ref{sqt:thm:density}.

\subsection*{[write] The uniformity of the imported localization}
\begin{writenote}[Added 6 October 2026, batch~108]
The uniform minor-arc bound~\eqref{sqt:eq:densityminor} is not a theorem
stated by Canfield and McKay: their Theorem~3 is an asymptotic statement for
$m,n\to\infty$ along admissible parameters, and the source supports
uniformity on $I$ by its own reading of their proof (``One can also check
uniformity directly in their localization proof'', the paragraph
before~\eqref{sqt:eq:densityminor}). Lemma~\ref{sqt:lem:sequential} derives
from the \emph{statement} of their Theorem~3 the weaker uniform bound that
the proof of Theorem~\ref{sqt:thm:density} actually uses;
Remark~\ref{sqt:rem:uniformity} reports what the write found in their proof.
\end{writenote}

\begin{lemma}[{[write]} Uniformity from the sequential statement]\label{sqt:lem:sequential}
Fix a compact $I=[\lambda_-,\lambda_+]\subset(0,\infty)$ and $\eta>0$. For
integers $n,s\ge1$ with $\lambda=s/n\in I$ let $\mathcal R_{\lambda,\eta}$ be
Canfield--McKay's region~(3.1) of \cite{CM} at $m=n$, $s=t$, with their
$\varepsilon$ equal to $\eta$, and put
\[
 E_\eta(n,s)=\frac1{I_{0,\lambda}}\int_{\mathcal R_{\lambda,\eta}^c}|\mathcal F_\lambda| .
\]
Then for every $p>0$,
$\sup\{n^pE_\eta(n,s):s\in\mathbb Z_{\ge1},\ s/n\in I\}\to0$ as $n\to\infty$.
The only input is \cite[Theorem~3]{CM}, read as a statement about every
sequence of admissible tuples with $m,n\to\infty$ along which~(1.1) holds and
$\lambda=n^{O(1)}$.
\end{lemma}
\begin{proof}
Suppose not. Then there are $p>0$, $c>0$, integers $n_1<n_2<\cdots$ and
$s_k$ with $s_k/n_k\in I$ and $n_k^pE_\eta(n_k,s_k)\ge c$ for all $k$.
Consider the tuples $(m,s;n,t)=(n_k,s_k;n_k,s_k)$; then $ms=nt$, and
$s=s(m,n)$ is a well-defined function on the pairs $(n_k,n_k)$ because the
$n_k$ are distinct. At $m=n$ the left side of~(1.1) is
$\tfrac23(4+1/u)$ with $u=\lambda(1+\lambda)\ge u_-=\lambda_-(1+\lambda_-)>0$,
so~(1.1) holds with $a=b=1/8$ as soon as
$\log n_k\ge\tfrac{16}3(4+1/u_-)$, and $\lambda_k\le\lambda_+$ gives
$\lambda=n^{O(1)}$. At $m=n$, $s=t$ the integrand $F$ of~(2.3) of
\cite{CM} is $\mathcal F_\lambda$ and their $I_0$ of~(4.1) is $I_{0,\lambda}$
(Section~\ref{sqt:sec:provenance}). Hence Theorem~3 of \cite{CM} applies along
this sequence: there is $\varepsilon'\in(0,\eta]$ (any sufficiently small one
will do) and a constant $C$ with
$\int_{\mathcal R_{\lambda_k,\varepsilon'}^c}|\mathcal F_{\lambda_k}|\le
Ce^{-n_k^{\varepsilon'}}I_{0,\lambda_k}$ for all large $k$. The three bounds
defining~(3.1) at $m=n$, namely $(1+\lambda)^{-1}n^{-1+4\varepsilon}$ and
twice $(1+\lambda)^{-1}n^{-1/2+\varepsilon}$, do not decrease as
$\varepsilon$ increases ($n\ge1$), and the half-circle condition does not
involve $\varepsilon$, so
$\mathcal R_{\lambda,\varepsilon'}\subseteq\mathcal R_{\lambda,\eta}$ and
$E_\eta(n_k,s_k)\le Ce^{-n_k^{\varepsilon'}}$. Then
$n_k^pE_\eta(n_k,s_k)\to0$, a contradiction.
\end{proof}

The proof of Theorem~\ref{sqt:thm:density} uses~\eqref{sqt:eq:densityminor}
only through the transfer Lemma~\ref{sqt:lem:transfer} with the scaled
weight, where the minor-arc integral enters twice (on
$\mathcal B_n\setminus\widetilde D$, after division by the gauge factor, and
outside the region) and is compared with $Q_{n,\lambda}Z_{n,\lambda}$, whose
ratio to $I_{0,\lambda}$ is bounded above and below on~$I$. Every later error
is $O_{I,p}(n^{-p})$, and the main term is of exact order $Z_{n,\lambda}$
because the exponent $H$ is real and uniformly bounded. So
Lemma~\ref{sqt:lem:sequential}, with $\eta=\varepsilon/2$, may replace the
uniform rate $O_I(e^{-c_In^\eta})$ in~\eqref{sqt:eq:densityminor}: the
torus-to-cube transfer then holds uniformly on $I$ with an error
$O_{I,p}(n^{-p})Z_{n,\lambda}$ for every $p$, and
Theorem~\ref{sqt:thm:density} and Corollary~\ref{sqt:cor:delta} follow as
printed. The same argument, with $\lambda=1$ fixed, is not needed for
Theorem~\ref{sqt:thm:main}, where Theorem~3 of \cite{CM} applies directly.

\begin{remark}[{[write]} What Canfield and McKay state, and what the write found in their proof]\label{sqt:rem:uniformity}
Theorem~3 of \cite{CM} assumes that ``$m,n\to\infty$ in such a way that~(1.1)
holds and $\lambda=n^{O(1)}$'' and concludes, ``for sufficiently small
$\varepsilon>0$'', that $\int_{\mathcal R^c}|F|=O(e^{-n^\varepsilon})I_0$;
their convention (p.~4) lets the implied constants depend on $a$, $b$
and~$\varepsilon$. No version uniform in a parameter family is stated, which
is why Lemma~\ref{sqt:lem:sequential} is needed. The write read their proof
of Theorem~3 (pp.~15--20 of the preprint). It bounds three pieces: the
region where at least $\tfrac13\min(mn^\varepsilon,m^\varepsilon n)$ pairs
have $\cos(\theta_j+\phi_k)\le\cos\delta$, by
$(2\pi)^{m+n}\exp\{-c_1(\log n)^{-1}\min(m^\varepsilon n,mn^\varepsilon)\}$
against $I_0=\exp\{O(m\log n+n\log m)\}$; the regions with $m_2+n_2\ge1$
exceptional coordinates, by a sum of
$n^{O(m_2+n_2)}\exp\{-c_3\lambda(1+\lambda)^{-1}(mn_2+nm_2)\}I_0$, giving
$O(e^{-n^{1-\varepsilon}})I_0$; and the rest of the complement of~(3.1), by
$e^{-n^\varepsilon}I_0$ through a factor
$\exp\{-\tfrac15\lambda(1+\lambda)^{-1}\min(m^{2\varepsilon},n^{2\varepsilon})\}$.
These are the three bounds the source displays before~\eqref{sqt:eq:densityminor}.
The parameter $\lambda$ enters only through $\lambda(1+\lambda)^{-1}$
(bounded below by $(\log n)^{-1}$ by~(1.1)), $A=O(\log n)$,
$N=\lceil6000(1+\lambda)\rceil$, $\delta=2\pi/N$ and their Lemmas~3 and~4,
of which Lemma~4 is stated ``uniformly for $\lambda>0$'' and both are stated
with proofs omitted (``We begin with two technical lemmas whose proofs are
omitted''). On a compact $I$ each of these quantities is bounded above and
below, so on the steps it read the write agrees with the source's reading;
but the implied constants (the $O(\cdot)$ exponents in
$n^{O(m_2+n_2)}$, the factor $O(\delta)$, the constants of Lemma~4) are not
tracked in their text, so the uniform \emph{rate}
$O_I(e^{-c_In^\eta})$ of~\eqref{sqt:eq:densityminor} remains the source's
reading of their proof, not a quoted theorem
(Question~\ref{sqt:q:uniform}). Nothing in Theorem~\ref{sqt:thm:density}
depends on that rate (Lemma~\ref{sqt:lem:sequential}). The same holds for
the source's ``The proof of Lemma~\ref{sqt:lem:transfer} now applies
verbatim with this weight'': it is the source's adaptation, whose two
normalization identities the write rechecked.
\end{remark}
